\section{Discussion}
\label{sec:discussion}

\subsection{Numerical precision is not physical identification}

A fitted edge may be numerically stable while remaining physically ambiguous. In the present spectra, coordinate perturbations move fitted edges by less than 1~meV, but admissible fitted observation models differ by approximately 6--7~meV. Reporting only optimizer uncertainty would therefore understate the dominant model dependence. The relevant uncertainty object is the declared candidate envelope together with boundary flags and missing specimen metadata.

\subsection{Fixed thresholds answer an operational question}

A fixed absorption crossing can be useful for reproducible device or process comparisons when the threshold and calibration are held constant. It is not automatically the latent band gap. The large threshold envelopes in the present analysis show that an unspecified threshold can dominate a few-meV comparison between material relations. The changes in nominal closest comparator in Sec.~\ref{sec:ranking-result} should be interpreted as a sensitivity result, not as evidence that a high crossing is physically superior.

\subsection{Composition and carrier metadata are part of the measurement}

The sub-meV Seiler--Hansen separation cannot be interpreted independently of composition uncertainty. Carrier state is also missing for both specimens, so filling-related contributions cannot be separated from the optical extraction. These are not ancillary reporting omissions; they are identifiability limits. The same logic applies to vacancy-sensitive processing history or a measured vacancy proxy when the absorption edge is used as a material-gap observation.

\subsection{Scope of the material-model comparison}

The analysis does not establish a preferred universal HgCdTe gap equation. Published Seiler is the nominal fitted-edge winner in this two-spectrum application, while fixed thresholds change that ordering. The provisional Hansen--Pade relation is retained only as a constrained research comparator. The Laurenti relation is retained as a high-confidence analytical reconstruction, with original coefficient covariance and validity limits unresolved. Hansen remains the zero-temperature baseline used by the provisional relation.

\subsection{Experimental consequence}

Further progress requires paired measurements rather than additional unpaired edge tables. The most informative next experiment measures a calibrated magneto-optical gap estimate and an absorption-derived edge on the same specimens while independently controlling or measuring composition, carrier state, and a vacancy-sensitive state variable. The paired factorial design converts the qualitative confounding diagram into a falsifiable acquisition program.

The design also supplies a stop rule for data accumulation. Another absorption-only source without composition uncertainty, carrier state, and a vacancy proxy cannot identify the declared terms, even if it adds many nominal edge values. Independent evidence should be acquired only when it changes a controlling conclusion or closes one of those missing contrasts.